\chapter{唐初：建国、帝国秩序与边疆}
\label{ch:early-tang}
\section{本章纲要：恢复帝国不是恢复旧日}

\paragraph{本章的论证。}唐初的稳定并非单靠“贞观之治”的美名，而是新朝将隋遗下的行政尺度，同关陇军政网络、关中政治中心、户籍赋役、律令和跨边疆外交重新编接。建国、治法、城市和西北扩张是同一件事的不同侧面：朝廷只有把地方资源、官僚程序与军队供给接起来，才有能力在内地建立可预期的秩序，并在草原和绿洲世界行动。

\paragraph{资源与约束。}长安、洛阳与关中军政传统提供了组织中心，州县、簿籍、粮运和成熟的书写行政提供了可复制的手段，丝路和边疆联盟则带来人力、马匹与交换机会。它们同时受制于关中供给、地区经济差异、旧精英与新官僚的竞争，以及维持北方和西方军事承诺的高昂成本；法典的条文和朝廷的任命，不能自动消除这些差异。

\paragraph{历史问题与章节后果。}本章追问：唐怎样在避免隋末式过度动员的同时，建立一个能向远方投送兵力、法令和礼仪的帝国？这一追问使唐初不再只是几位明君的连续故事，而成为中央与地方、内地与边疆相互塑造的过程。其后果是，帝国秩序越能把资源送往边境和都城，继承、任官和财政的每一次调整便越会影响整个网络，为武周与开元时期的重组准备条件。

\paragraph{史料提示。}《旧唐书》《新唐书》与《资治通鉴》保存朝廷行动的主线，却常以成熟帝国或后来的政治判断整理唐初；律令和会要显示规范，不等于各地实践。墓志、敦煌吐鲁番文书、碑刻、考古材料和外来文本能显出地方与跨区域活动，但保存不均，不能替所有群体发言。本章据此保留材料之间的距离。

\section{从太原到洛阳：建国不是一日的受禅}
\label{sec:early-tang-foundation}

617年李渊自太原起兵，入关后拥立代王杨侑并受封唐王；次年受禅建唐。\eventanchor{ST-617-LI-YUAN-TAIYUAN}{太原起兵（617）}太原军镇、关陇网络和与突厥的交涉，使他能够取得长安，却不等于已经取得一个可直接调度的帝国。唐初史书特别突出李世民的作用和突厥关系，正是后来的继承与军功记忆塑造了材料；可以把握的是关中出现了一个能接收隋朝官署与财政名义的中心。%
\footnote{太原起兵与唐建国见 ST-617-LI-YUAN-TAIYUAN、ST-618-TANG-FOUNDING，依据《旧唐书·高祖本纪》《新唐书·高祖本纪》及《资治通鉴》隋、唐纪。结盟条款和个人作用的细节受唐代叙事重塑。}

\begin{event}{太原起兵与建唐：617--618年，关中}
起兵、入关和受禅构成可相互校读的年序；它们并不证明新朝已经在各地取得同样的资源、服从或行政能力。后文循官署、道路与军政网络追问的，正是这项转换如何被继续完成。
\end{event}

王世充在洛阳建郑，窦建德的夏军又能自河北南援；621年虎牢之战及王世充降唐，使唐取得洛阳与中原仓储，却没有终结地方接管。次年刘黑闼在河北再起，说明军队、旧部和地方征发的关系不能由一纸降书消除。唐朝的统一战争因此不是一路“平定”，而是围绕洛阳、河北道路与州县人事反复建立可用控制。%
\footnote{郑建国、虎牢与河北复叛见 ST-619-ZHENG-FOUNDING、ST-621-HULAO、ST-622-LIU-HEITA，依据《旧唐书》王世充、窦建德、刘黑闼诸传、《新唐书》相关列传及《资治通鉴》。兵数、突阵和地方支持度主要来自胜者的军功叙事。}

626年玄武门政变将继承争夺置于禁军和宫城的暴力中；两月后颉利可汗至渭水北岸，唐廷以会盟换得退兵。二者相接，正说明新皇帝不能把内部继承与边防分开处理。628年梁师都政权覆灭，才使关中北缘较为稳固；唐的早期国家能力是在军队、任官、粮储与边地外交都尚未定型时逐步取得的，而不是贞观一开始便自然完备。%
\footnote{玄武门、渭水与梁师都覆亡见 ST-626-XUANWU、ST-626-WEISHUI、ST-628-LIANG-SHIDU，依据《旧唐书·太宗本纪上》《旧唐书·突厥传上》《旧唐书·梁师都传》及《资治通鉴》。政变预谋和渭水谈判细节均有后出的修辞层次。}

\subsection{交叉阅读窗口二：受禅、律文与滕王阁的远望}

617--628年的建国叙事，宜以《旧唐书》《新唐书》高祖、太宗纪和《资治通鉴》的年序相互牵制：前者留下唐朝开国者及北宋重修者的不同取景，后者将兵变、受禅与边事压成因果链。永徽四年（653）颁定的《唐律疏议》与后出《唐会要》所辑制度材料，不能倒过来证明621年每一州县已经依同一格式裁判；它们却使我们能问，唐为何要把隋以来的官署、户婚和刑名转化为可复用的行政语言。初唐墓志所见籍贯、官衔与婚姻可作局部校验，但其作者、受葬者皆有选择，不能补成统一战争中普通家庭的档案。%
\footnote{建国与统一战争可互校《旧唐书·高祖本纪》《旧唐书·太宗本纪上》《新唐书》相关本纪、列传与《资治通鉴》唐纪；律文见《唐律疏议》，制度条目可参《唐会要》。永徽律的年代晚于本节多数事件，故这里只把它当作制度延续的提问，不当作617--628年的直接证词。}

王勃《滕王阁序》有“关山难越，谁悲失路之人；萍水相逢，尽是他乡之客”之句。此篇为骈文宴集之作，传统系于675年前后，作者以早慧而失意的宾客身份铺陈江山、仕途与相逢；它的句法把交通节点的壮阔转成个人去留的压力。它不能证明唐初某条道路的人口流向，时间也晚于虎牢和渭水数十年；但它提醒读者，统一不是把关山抹平，而是让道路、州县名分和仕进机会重新决定谁能够抵达中心。对钱穆所问中央制度如何获得地方支点而言，这是一种可检验的空间问题；对吕思勉的社会史问题而言，则须回到律文、墓志和文书，而不能让一篇华丽骈文代替迁徙史。

\section{法律、宗教与继承：秩序的不同语言}
\label{sec:tang-imperial-order}

贞观元年修定律令、永徽四年完成《唐律疏议》，使初唐官府拥有较稳定的刑名与解释框架。法典的连续性不能被误读为社会的整齐：律文说明规范、疏议说明官方解释，而州县诉讼、身份登记和实际裁断仍须由文书与个案来观察。唐朝对隋制的继承，最可见于这种把既有制度重新编排为可援引语言的工作。%
\footnote{贞观律与《唐律疏议》见 ST-627-ZHENGUAN-CODE、ST-653-TANGLU-SHUYI，依据《旧唐书·刑法志》《唐律疏议》及《通典·刑法》。颁行与地方执行不可混同。}

635年阿罗本到长安获准译经，664年玄奘卒于玉华宫而译场与弟子网络仍在延续。这两件事不宜被写成唐廷对所有宗教一概“开放”的证明：许可、译经、寺院、施主与官府之间各有时段和区域的条件。它们只表明都城及其交通网络能够容纳多种文本、僧侣与仪式进入政治可见的空间；景教碑和佛教传记又各有自我说明与典范化的目的。%
\footnote{阿罗本与玄奘见 ST-635-NESTORIAN、ST-664-XUANZANG-DEATH，依据《旧唐书》相关记载、景教碑及《大慈恩寺三藏法师传》《续高僧传》等。入华、获准与具体社群规模须分开讨论。}

643年李承乾被废、649年李治即位，则提醒人们法令与宗教并不能消除皇室继承的风险。传记把谋反案写得层次分明，却难以恢复宫廷各派全部打算；较可靠的是太宗通过废立重新指定继承人，高宗即位后又须依托既有官僚和法制语言取得承认。帝国秩序不是单一制度的结果，而是在法庭、寺院、宫门和人事任命中不断被重申。%
\footnote{承乾案与高宗即位见 ST-643-CHENGQIAN、ST-649-TAIZONG-DEATH，依据《旧唐书·太宗本纪》《旧唐书·承乾传》《新唐书》相关纪传及《资治通鉴》。案情与人物动机受继承结果影响，不能作心理断言。}

\section{胜利之后，边疆仍要供给}
\label{sec:tang-frontiers-silk-road}

630年唐军破东突厥，次年安置降附部众；634年击吐谷浑，638年松州之战后又转入与吐蕃的婚盟。这里没有一条从胜利到服从的直线：俘获可改变草原政治，迁置则立刻提出牧地、粮食、任官与身份分类的问题；军事对峙之后的婚盟，也不是一方永久支配另一方的证明。汉文正史多以唐廷为观察中心，突厥与吐蕃的行动逻辑须从碑铭、文书和外部材料中补问。%
\footnote{东突厥、安置、吐谷浑、松州与文成公主入吐蕃见 ST-630-EASTERN-TURKS、ST-631-TURK-SETTLEMENT、ST-634-TUYUHUN、ST-638-SONGZHOU-TIBET、ST-641-WENCHENG，依据《旧唐书》《新唐书》突厥、吐谷浑、吐蕃传与《资治通鉴》。人口、贡赐和婚盟的实际条件不宜照录为精确数字。}

高昌在640年覆亡并置西州，648年攻龟兹，657年又破西突厥；这些推进把都护、羁縻与地方州县并置在天山南北。直接置州、任都护或册立首领，分别意味着不同的行政深度，不能合称为一个无缝的“西域统治”。唐军645年久攻安市未克而撤出辽东，也提醒人们远征的极限：寒冷、粮道和城防足以抵消一时的兵力优势。%
\footnote{高昌、辽东、龟兹与西突厥事件见 ST-640-GAOCHANG、ST-644-GOGURYEO-EXPEDITION、ST-645-LIAODONG-WITHDRAWAL、ST-648-KUCHA、ST-657-WESTERN-TURKS，依据《旧唐书》相关列传、《新唐书》地理及外族传和《资治通鉴》。都护辖境、兵数与战果具有行政和军功修辞。}

在东北，唐与新罗的联合先后灭百济、在白江口击败倭援军、再攻破平壤。联军的年序可以确证，却不能把新罗、倭和百济复国军写成唐朝叙事中的附属角色；《三国史记》《日本书纪》同样有本国的编纂立场。668年以后安东都护府的设置也不是高句丽社会被立刻同质接收的凭证，而是战争胜利转化为行政安排时产生的新摩擦。%
\footnote{百济、白江口与高句丽覆亡见 ST-660-BAEKJE、ST-663-BAEKGANG、ST-668-GOGURYEO-CONQUEST，依据两《唐书》东夷传、《三国史记》《日本书纪》及《资治通鉴》。战场人数、航路和地方反应在不同传统中有异。}

670年大非川战败、676年前后河湟后撤，679年东突厥复兴；武周在692年恢复安西四镇、702年置北庭，而营州契丹起事又在695--697年暴露东北的压力。边疆不是一面不断向外推的墙，而是若干需分别供给和谈判的走廊。747年、751年对南诏的两次失利，以及怛罗斯之败，同样不可被压缩为帝国“由盛转衰”的单一信号：云南山地、河中交通、盟友关系和远征补给的限制各不相同，却共同限制了中央把军令变成稳定控制的能力。%
\footnote{吐蕃、后突厥、安西北庭、契丹、南诏与怛罗斯诸事见 ST-670-DAFEICHUAN、ST-676-HEHUANG-LOSS、ST-679-ASHIDE-REBELLION、ST-692-ANXI-RECOVERY、ST-695-KHITAN-REBELLION、ST-697-KHITAN-SUPPRESSION、ST-702-BEITING、ST-747-NANZHAO、ST-751-TALAS、ST-751-NANZHAO-DEFEAT。核心为两《唐书》诸传与《资治通鉴》；吐蕃、突厥、回鹘碑铭及中亚材料是必要校正，不能以唐方失败或胜利叙事代替他方目的。}

\phantomsection\label{fig:vol03:early-tang:tang-frontiers-silk-road}
\begin{center}
\begin{tikzpicture}[
  x=.78cm,y=.78cm,font=\small,
  place/.style={draw=timeline!85,fill=paper,rounded corners=2pt,align=center,
    inner sep=3pt},
  court/.style={place,draw=dynasty,fill=dynasty!22},
  frontier/.style={place,fill=timeline!18},
  external/.style={place,draw=source,fill=source!13},
  note/.style={font=\scriptsize,text=source,align=center}
]
  \fill[paper] (0,0) rectangle (12.5,6.9);
  \node[font=\bfseries\large,text=dynasty] at (6.25,6.42)
    {唐初的边疆走廊与跨区域战争};
  \node[note] at (6.25,5.94) {630--668年若干军事、行政与外交节点（示意）};

  \fill[source!11] (.55,4.73) -- (11.95,4.55) -- (11.8,5.6) -- (.72,5.76) -- cycle;
  \node[note] at (7.65,5.15) {草原、天山与东北方向：不同政治行动者和道路的空间};

  \node[court] (changan) at (4.85,2.85) {长安\\朝廷与军政调度};
  \node[frontier] (lingzhou) at (4.1,4.52) {北方边州\\东突厥战争后};
  \node[external] (turks) at (5.95,5.0) {东突厥\\部众、牧地与交涉};
  \node[frontier] (songzhou) at (2.48,2.3) {松州、河湟\\吐谷浑与吐蕃方向};
  \node[external] (tibet) at (.86,1.48) {吐蕃方向\\战争、婚盟与道路};
  \node[frontier] (xizhou) at (7.0,3.82) {高昌／西州\\640年置州};
  \node[frontier] (kucha) at (9.15,3.58) {龟兹、安西\\648年后};
  \node[external] (western) at (10.92,4.78) {西突厥方向\\657年战争与重组};
  \node[frontier] (liaodong) at (8.36,1.42) {辽东、平壤\\远征与战争};
  \node[external] (silla) at (10.78,1.1) {百济、新罗、倭\\联盟与海上战争};

  \draw[<->,thick,dashed,timeline] (lingzhou) -- node[above,sloped,note]
    {战争后的安置与交涉} (turks);
  \draw[->,thick,dynasty] (changan) -- node[left,note]
    {北方守御与任官} (lingzhou);
  \draw[<->,thick,dashed,source] (songzhou) -- node[below,sloped,note]
    {638年后：战争、婚盟与边地协商} (tibet);
  \draw[->,thick,dynasty] (changan) -- node[above,sloped,note]
    {640年高昌覆亡、置西州} (xizhou);
  \draw[->,thick,dynasty] (xizhou) -- node[above,sloped,note]
    {648年龟兹战争与驻防} (kucha);
  \draw[<->,thick,dashed,timeline] (kucha) -- node[above,sloped,note]
    {天山南北的多重关系} (western);
  \draw[->,thick,dynasty] (changan) -- node[below,sloped,note]
    {645年辽东远征（约略）} (liaodong);
  \draw[<->,thick,dashed,source] (liaodong) -- node[below,sloped,note]
    {660--668年联盟与战事} (silla);

  \node[draw=source!75,fill=paper,rounded corners=2pt,align=center,
    inner sep=4pt,font=\scriptsize,text=source] at (2.5,5.03)
    {节点与箭头只提示可见关系；\\不表示稳定统属、精确路线或连续疆界};
\end{tikzpicture}

\smallskip
\small\textit{图\quad 唐初的边疆走廊与跨区域战争（630--668，示意）：据《旧唐书》《新唐书》突厥、吐谷浑、吐蕃、西域及东夷相关列传和《资治通鉴》，并参照突厥碑铭、吐蕃材料、《三国史记》与《日本书纪》，标出东突厥、高昌、西州、龟兹、辽东及朝鲜半岛方向的若干事件节点。官修唐史以朝廷行动为轴，外部材料各有自己的编年和政治语言；故本图不绘制唐的精确疆界或都护辖境，也不将战役、置州、联盟或婚盟等同于稳定控制、单向服从或完整行军路线。}
\end{center}

\section{初唐的共同语言与不均等的经验}
\label{sec:early-tang-society-institutions}

贞观与永徽时期的律令、官制和都城常被作为一个“成熟帝国”的证据。更接近材料的读法，是把它们看作使不同人能够同官府打交道的共同语言：户婚、刑名、官司与仓库有了可援引的分类，官员也能以同一套文书解释职责。可援引不等于可抵达。唐律疏议保存的是规范及其官方解释，敦煌、吐鲁番文书保存的是特定地点的实践；二者之间的距离，正是地方家庭、书吏和官员变通制度的空间。%
\footnote{贞观律令与《唐律疏议》见 ST-627-ZHENGUAN-CODE、ST-653-TANGLU-SHUYI，核心材料为《旧唐书·刑法志》《唐律疏议》《通典》及令典残卷。不能从成文法直接推出全国司法实践。}

长安也不是只由皇帝和宰相居住的舞台。635年阿罗本获准译经、664年玄奘死后译场继续活动，说明僧侣、译者、施主、官府和道路把不同语言的宗教知识接到都城；它们并不证明国家为所有社群提供同等空间。景教碑、僧传和寺院题记都是各自的自我表述，适于显示网络与仪式，不能据此计算信徒人数，或把一座都城的情况外推为整个乡村社会。%
\footnote{阿罗本与玄奘的节点见 ST-635-NESTORIAN、ST-664-XUANZANG-DEATH，依据《旧唐书》相关记载、景教碑、《大慈恩寺三藏法师传》及《续高僧传》。文本成书、碑刻立置和实际活动的日期需要分别辨认。}

皇室继承同样穿过这些制度，而非停留在制度之外。李承乾被废、李治继位，后来武后进入皇后与临朝位置，表明法典、礼仪和官僚网络可以提供争论的格式，却不能预先消除继承中的暴力与排除。初唐国家最值得追问的不是它是否已经“强大”，而是它如何把法律、城市宗教、家族地位和边防军功暂时接在一起；这个接合让帝国能够行动，也使被排除者和地方差异继续存在于文本之外。%
\footnote{承乾案、高宗即位及武后立后见 ST-643-CHENGQIAN、ST-649-TAIZONG-DEATH、ST-655-WU-EMPRESS，依据两《唐书》本纪、列传及《资治通鉴》。宫廷案件由后来的胜者保存，动机和言辞不能作确定复原。}


\section{章末综合：可预期的秩序与扩张的价格}

\paragraph{制度性的综合。}唐初的律令、官制和地方行政不是静止的“制度遗产”，而是让户婚、刑名、赋役、任官和军需能够被分类、传递并在争议中援引的共同语言。它降低了重新统一后的协调成本，却仍须依赖地方书吏、家族、仓场和军镇的实际合作；越是广泛的边疆行动，越会检验这种合作能否被持续供给。

\paragraph{历史性的综合。}因此，唐初的成就不是已经解决了帝国的规模问题，而是把问题置于一套更稳固也更昂贵的组织之中。下一章将考察皇位、礼仪、人才与财政如何在武周和开元被重新配置，并说明这种弹性何以同时扩大了世界帝国的行动范围。
